\chapter{Phân công và tỷ lệ đóng góp}
\label{app:phan-cong}

Bảng dưới ghi phân công chính của nhóm ba thành viên. Tỷ lệ đóng góp chưa chốt tại thời điểm
viết; nhóm sẽ tính từ lịch sử commit và xác nhận trước khi nộp. Bảng này để người đọc biết ai
phụ trách phần nào, không dùng để so sánh mức đóng góp.

\begin{table}[H]
  \centering\small
  \begin{tabularx}{\textwidth}{@{}>{\raggedright\arraybackslash}p{.22\textwidth} >{\raggedright\arraybackslash}X >{\raggedright\arraybackslash}X@{}}
    \toprule
    \textbf{Thành viên} & \textbf{Phụ trách chính} & \textbf{Phụ trách phụ} \\
    \midrule
    Huỳnh Thanh Tiến & Engine bitboard, interface tác tử, CI, UI; cấp 4 (MCTS), cấp 7--8 (PUCT, self-play) & README, demo; A3/A4 \\
    Võ Thị Như Ý & Cấp 2, cấp 3 (Alpha-Beta), A1/A2 & Phân tích lỗi tác tử \\
    Đặng Ngọc Như Ý & Cấp 5--6, arena và thống kê & Notebook, biểu đồ \\
    \bottomrule
  \end{tabularx}
  \caption{Phân công chính của nhóm. Tỷ lệ phần trăm sẽ được chốt từ lịch sử commit khi nộp.}
  \label{tab:phan-cong}
\end{table}

\section{Cách xác định tỷ lệ đóng góp}
Nhóm dùng ba nguồn: (i) số commit và số dòng thay đổi theo từng thư mục; (ii) issue được giao
và issue đã đóng; (iii) xác nhận chéo của các thành viên. Không dùng số dòng tuyệt đối làm
trọng số duy nhất, vì một thay đổi nhỏ ở engine có thể quan trọng hơn nhiều dòng notebook. Tỷ
lệ cuối cùng do nhóm thống nhất và ghi ở bản nộp.

\section{Bằng chứng đóng góp}
Lịch sử Git trong repo là bằng chứng chính. Mỗi commit ghi người thực hiện và thư mục thay đổi;
các nhãn issue (TV1--TV4) trong đề bài và trong repo giúp đối chiếu. Báo cáo không ghi tỷ lệ
ước lượng khi chưa có xác nhận, để tránh con số không kiểm chứng được.